\section{Discussion}
\label{sec:discussion}
This chapter interprets the DTSE results by linking the stress-response patterns from Chapter~8 to mechanism design differences between BME-oriented and capped-supply profiles. The aim is not to attach simple labels such as ``robust'' or ``fragile.'' Instead, the discussion traces how stress moves through each mechanism: which metric families react first, how those deviations spread, and which operational failure-mode signatures emerge under the DTSE contract.

The discussion stays inside the same interpretive boundaries used throughout the thesis. Mechanism facts remain grounded in documentation, modeled assumptions remain explicit, and DTSE outputs remain conditional on the frozen scenario grid and thresholds. Within those boundaries, the chapter asks what the observed stress signatures imply for DePIN robustness, governance response, and the practical use of DTSE as an evaluator.

\subsection{Interpretive Contract and Claim Classes}
This discussion builds on the bridge established in Section~\ref{subsec:deferred_interpretation_register}. To avoid category errors, statements are kept within three claim classes:
\begin{itemize}
    \item \emph{Mechanism facts}: protocol-defined rules (issuance, reward allocation, and sinks) supported by protocol documentation, not by DTSE output \cite{OnocoyWhitepaper301, OnocoyToken, OnocoyRewards}.
    \item \emph{Modeled assumptions}: abstractions introduced for experimental control (e.g., demand regimes, reduced-form price signals, provider-cost distributions, and decision rules), not claims about real-world prevalence \cite{Morris2019, Braakman2022}.
    \item \emph{DTSE outputs}: descriptive patterns in state variables and evaluation metrics produced under the frozen scenario inputs, used for comparative sensitivity statements inside the DTSE experiment set.
\end{itemize}
Accordingly, interpretations remain bounded by the stress-channel definitions (Section~\ref{sec:stress_scenarios}) and metric definitions (Section~\ref{sec:evaluation_metrics}). They are discussion-level readings of the experiment set, not empirical measurements of live-network behavior.

\subsection{Human Intervention Archetypes Under Stress}
Chapter~8 formalized failure modes as operational signatures: recurring metric patterns under specified stress inputs (Table~\ref{tab:failure_modes}). This subsection introduces \emph{human intervention archetypes}: decision-relevant response patterns for discussing what governance might try as a system approaches those signatures. They are analytical categories used to organize the discussion, not claims about universal real-world behavior.

Because DTSE does not model intent or off-chain governance constraints, the archetypes below should be read as response classes that become representable in DTSE only when translated into explicit rule changes or exogenous-input changes. For ONO-specific interpretation in this section, ``emissions'' refer to modeled reward distributions from a fixed inventory path rather than uncapped new token creation.

\subsubsection*{Subsidy Continuation}
\emph{Definition.} Governance leaves the incentive rule set unchanged while usage weakens.

\emph{DTSE anchor.} This posture maps to \textbf{Reward--Demand Decoupling} (Table~\ref{tab:failure_modes}): utilization and usage-linked sink fields fall versus baseline while distribution remains active under mechanism rules, followed by deterioration in incentive--demand alignment proxies and margin compression in provider profitability. In the demand-contraction channel, ONO registers stress first in utilization and sinks, with participation effects (active provider count) appearing materially later (Section~\ref{stress-scenario-results-comparative-outcomes}; Table~\ref{tab:dtse_signal_timing_ono}).

\emph{Interpretation.} This label names a modeled stance rather than an intention. In practice, delayed adjustment may result from governance latency or treasury constraints, but those frictions remain contextual explanation rather than load-bearing evidence in this thesis.

\subsubsection*{Broad Incentive Increase}
\emph{Definition.} Governance increases aggregate incentive intensity to counter participation erosion under stress.

\emph{DTSE anchor.} The run set does not simulate counterfactual policy interventions, so this archetype is interpreted as a response \emph{to} the observed signatures rather than as a simulated mechanism change. The relevant evidence is that the competitive-yield and cost-inflation channels register first in provider-economics and participation fields (provider profit, provider churn, active provider count) under the frozen thresholds (Table~\ref{tab:dtse_cross_scenario_synthesis}; Table~\ref{tab:dtse_signal_timing_ono}). These are the same field families that underwrite \textbf{Profitability-Induced Churn} and \textbf{Elastic Provider Exit} (Table~\ref{tab:failure_modes}).

\emph{Interpretation.} A broad incentive increase targets the first-moving fields in these channels, but it can also weaken incentive--demand alignment if reward outflow expands while usage-linked sinks do not recover (Section~\ref{stress-scenario-results-comparative-outcomes}). Outside DTSE, short-horizon participation support can conflict with budget discipline and governance risk tolerance. Formal counterfactual intervention modeling therefore remains a future extension (Section~\ref{subsec:future_research_agenda}).

\subsubsection*{Incentive Re-Targeting}
\emph{Definition.} Governance changes reward allocation rather than increasing aggregate reward magnitude, for example by concentrating incentives on higher-commitment capacity or validated service outcomes.

\emph{DTSE anchor.} The DTSE experiment set does not simulate allocation-policy counterfactuals, so this archetype is again read as a response \emph{to} observed participation-side signatures. The motivating evidence is that competitive-yield pressure registers primarily in participation fields rather than as abrupt service discontinuities (Table~\ref{tab:dtse_signal_timing_ono}; Table~\ref{tab:dtse_cross_scenario_synthesis}). These same fields underwrite \textbf{Elastic Provider Exit} and, under time-order conditions, \textbf{Latent Capacity Degradation} (Table~\ref{tab:failure_modes}).

\emph{Interpretation.} Re-targeting aims to reduce exit elasticity by shifting rewards toward providers whose continuation improves redundancy or verified service output. The effect is to change the distribution of profitability and churn rather than the mean reward level. Within DTSE, such an action becomes representable only if it is translated into explicit changes to allocation rules, thresholds, or provider-type response parameters. Related mechanism-design considerations around targeted incentives and participation dynamics are discussed in the literature \cite{Tan2020, Chiu2024}. Outside DTSE, implementation complexity---including verification criteria, monitoring quality, and communication overhead---can constrain how quickly re-targeting can be deployed.

\subsubsection*{Non-Structural Response}
\emph{Definition.} Governance relies on actions that do not change mechanism rules or exogenous conditions, such as communication strategies, signaling, or informal coordination without parameter change.

\emph{DTSE anchor.} DTSE does not model these actions directly. Within the simulator, the archetype is therefore equivalent to no modeled intervention: the system continues under the same stress inputs, and the output patterns cannot evaluate communication or signaling unless those actions are translated into modeled parameter, exogenous-input, or decision-rule changes.

\emph{Interpretation.} Non-structural responses may delay behavioral shifts in practice, but DTSE cannot attribute changes to communication alone. A practical friction outside DTSE is credibility decay: communication-only responses can lose effect quickly if operators continue to observe deteriorating margins or rising churn signals.

\subsubsection*{Coordinated Operational Intervention}
\emph{Definition.} Governance coordinates actions that change operational conditions treated as exogenous drivers in DTSE, such as provider cost structures, effective capacity, or service-demand trajectories.

\emph{DTSE anchor.} In this thesis, operational interventions correspond to counterfactual adjustments of exogenous inputs (demand regimes, cost inflation, liquidity depth assumptions) rather than mechanism rule changes. The results chapter shows demand-side stress registering first in utilization-linked channels and cost-side stress registering first in provider profitability (Section~\ref{stress-scenario-results-comparative-outcomes}; Table~\ref{tab:dtse_cross_scenario_synthesis}), motivating an archetype framed in demand and cost terms. DTSE can represent operational interventions only when they map directly into scenario-grid input changes.

\emph{Interpretation.} This archetype clarifies the boundary between what DTSE can encode as an intervention, namely a changed exogenous trajectory, and what remains out of scope, namely governance dynamics that do not alter modeled inputs. Outside DTSE, a primary friction is coordination capacity: these interventions may require simultaneous changes across partners, operations, and treasury policy.

\medskip
Taken together, these archetypes provide a neutral vocabulary for discussing responses to DTSE-identified stress signatures without intent labels or empirical prevalence claims. They also separate intervention types that are representable in DTSE (rule changes and exogenous-input changes, when explicitly parameterized) from actions that remain out of scope (unmodeled governance dynamics and off-chain behavior).

\subsection{Onocoy-Specific Interpretive Signals}
Onocoy serves as the empirical anchor for mechanism mapping and for interpreting stress-test relevance in a hardware-dependent DePIN setting. This subsection introduces no new mechanism claims. Instead, it clarifies how the DTSE metric vocabulary is applied to a capped-supply profile when the underlying mechanism facts are sourced from official documentation (Section~\ref{sec:onocoy_case}).

In this thesis, an \emph{Onocoy-specific interpretive signal} is an observable quantity (or a documented rule parameter) that can be mapped into a DTSE metric family without requiring behavioral inference. The mapping is organized into three classes:
\begin{enumerate}
    \item \textbf{Usage-linked sinks/credits:} protocol-defined usage or service credits that correspond to usage-linked sinks in DTSE (burn-linked measures and incentive-solvency proxy families).
    \item \textbf{Reward rules:} documented distribution logic that corresponds to DTSE ``emissions'' and net-emission dynamics. For ONO, ``emissions'' refer to distributions from a fixed inventory path rather than uncapped issuance.
    \item \textbf{Eligibility/verification requirements:} documented participation conditions that map to participation thresholds or effective-capacity contributions in DTSE. Where the mapping requires values not specified in documentation, parameters are treated explicitly as modeled assumptions.
\end{enumerate}

With this mapping in place, the operational diagnostic matrix (Table~\ref{tab:failure_modes}) becomes a practical lens for reading DTSE outputs in ONO terms:
\begin{itemize}
    \item \textbf{Reward--Demand Decoupling:} usage-linked fields (utilization and sink proxies) weaken versus baseline while reward distributions remain active, indicating a deterioration in incentive--usage alignment under the modeled channel.
    \item \textbf{Profitability-Induced Churn:} provider profit remains structurally negative for sustained windows under the cost and price-process assumptions, with participation stress registering through provider churn and active provider count.
    \item \textbf{Liquidity-Driven Compression:} a modeled price dislocation compresses token-denominated reward value and propagates into provider-economics and participation fields in the same or adjacent windows under frozen thresholds.
    \item \textbf{Elastic Provider Exit:} participation responds sharply to outside-opportunity pressure under the modeled switching rule, with churn and active provider count triggering before utilization becomes the first-moving field.
    \item \textbf{Latent Capacity Degradation:} participation erosion appears first while service fields deteriorate later, consistent with redundancy being consumed before visible capacity shortfalls emerge.
\end{itemize}

These signals define how mechanism facts are instantiated in DTSE and interpreted within the model contract. Where the mapping requires parameter values not specified in official documentation, those values are treated as modeled assumptions with explicit bounds and sensitivity handling. Assumption bounds and sensitivity intervals are listed in Appendix Section~\ref{app:dtse_config_register}. Concrete ONO-to-DTSE mapping instances (signal to DTSE metric family) are provided in Section~\ref{sec:onocoy_case} and Appendix Sections~\ref{app:dtse_config_register} and~\ref{app:dtse_override_ledger}.

\subsection{External Validity, Generalization Limits, and Evidence Boundaries}
DTSE is a comparative evaluator: it tests how mechanism profiles respond to standardized stress channels under controlled assumptions. External validity is therefore bounded by construction. The cross-scenario synthesis (Table~\ref{tab:dtse_cross_scenario_synthesis}) summarizes metric-family sensitivity and detection timing \emph{inside} the DTSE experiment set; it does not imply that the same ordering must hold in live systems.

The main generalization limits are:
\begin{enumerate}
    \item \textbf{Demand abstraction:} demand regimes are exogenous stylizations; DTSE does not endogenize user formation, price elasticity, or adoption feedback beyond explicit parameterization \cite{Morris2019, Braakman2022, Collins2024}.
    \item \textbf{Price-process abstraction:} the modeled price series transmits stress into provider margins but does not represent market microstructure, strategic trading, or institutionally mediated liquidity in detail \cite{Morris2019, Braakman2022, Collins2024}.
    \item \textbf{Behavioral simplification:} provider behavior is represented through reduced-form decision rules; coordination and multi-homing enter only through modeled thresholds and switching logic \cite{Morris2019, Braakman2022, Collins2024}.
    \item \textbf{Governance scope limits:} governance and verification dynamics remain out of scope unless explicitly represented as DTSE inputs and supported by protocol documentation \cite{Morris2019, Braakman2022, Collins2024}.
\end{enumerate}

Two boundaries are specific to the Onocoy anchor case. First, DTSE evaluates tokenomic resilience, but it does not simulate the cyber-physical verification layer, such as anti-spoofing controls or data-plane routing, which can add operational friction in deployment. Second, where public documentation leaves mechanism rules underspecified, such as staking or collateral granularity, those elements remain structural unknowns unless introduced as explicit bounded assumptions. Threshold outputs are therefore modeled sensitivities under stated assumptions rather than empirical guarantees \cite{OnocoyWhitepaper301, OnocoyRewards, Chiu2024}.

Table~\ref{tab:limitation_claim_impact_map} maps these limitations to affected claim classes and the corresponding audit paths used in this thesis.

\begin{table}[ht]
    \centering
    \footnotesize
    \renewcommand{\arraystretch}{1.2}
    \setlength{\tabcolsep}{4pt}
    \begin{tabularx}{\textwidth}{@{}>{\raggedright\arraybackslash}p{0.19\textwidth}>{\raggedright\arraybackslash}p{0.19\textwidth}>{\raggedright\arraybackslash}p{0.25\textwidth}X@{}}
        \toprule
        \textbf{Limitation} & \textbf{Affected Claim Class} & \textbf{Interpretive Impact} & \textbf{Mitigation / Audit Path} \\
        \midrule
        Exogenous demand stylization &
        Modeled assumptions; DTSE outputs &
        Limits claims about endogenous adoption dynamics and demand elasticity outside modeled scenarios &
        Keep claims baseline-relative and scenario-conditional; audit scenario inputs in Appendix Section~\ref{app:dtse_scenario_grid}. \\
        \addlinespace
        Reduced-form price process &
        Modeled assumptions; DTSE outputs &
        Limits claims about market microstructure and strategic trading behavior &
        Restrict interpretation to incentive-transmission patterns; audit formulas and thresholds in Appendix Sections~\ref{app:dtse_scenario_formula_map} and~\ref{app:dtse_threshold_table}. \\
        \addlinespace
        Simplified provider decision rules &
        Modeled assumptions; DTSE outputs &
        Limits behavioral realism for coordination, multi-homing, and strategic adaptation &
        Use directional pattern language; audit profile assumptions and overrides in Appendix Sections~\ref{app:dtse_config_register} and~\ref{app:dtse_override_ledger}. \\
        \addlinespace
        Governance and verification out of scope &
        Mechanism facts; modeled assumptions; DTSE outputs &
        Prevents causal interpretation of governance-response efficacy from DTSE results alone &
        Keep governance statements as scoped interpretation; route intervention modeling extensions to Section~\ref{subsec:future_research_agenda}. \\
        \bottomrule
    \end{tabularx}
    \caption{Limitation-to-claim impact map for discussion-level interpretation.}
    \label{tab:limitation_claim_impact_map}
\end{table}

Two interpretation constraints from Chapter~8 remain decisive when generalizing results: (i) baseline trajectories can drift under neutral inputs, so scenarios are read as deviation patterns rather than stable-versus-unstable classifications; and (ii) selected cross-profile metrics can occupy different numeric scales in specific channels, so cross-profile statements emphasize directional patterns under matched inputs rather than scalar score ordering.

These boundaries shape how the thesis answers the main question. The thesis does not claim universal superiority of any mechanism class. It uses DTSE to compare BME-oriented and capped-supply profiles under standardized stress channels and summarizes behavior through operational failure-mode signatures. Conclusions are therefore comparative statements about relative robustness \emph{within the experiment set}, conditional on the scenario grid and parameter ranges used.

\subsubsection*{Evaluator Usage Stance}
DTSE is used here as a decision-support evaluator: run matched stress channels, identify the earliest-trigger metric families and associated failure-mode signatures, and localize where sensitivity concentrates before design or governance changes are considered. This stance preserves claim-class discipline and artifact traceability while avoiding interpretation beyond modeled assumptions.

\subsection{Reproducibility and Artifact Traceability}
DTSE's main methodological strength is traceability. The methodology fixes parameterization, scenario inputs, and random-seed policy, and it reports outcomes using distributional summaries rather than single-run narratives. Each interpretive claim in this chapter is therefore anchored to a scenario identifier, a metric definition and computation rule, and an output artifact that displays the relevant pattern. The audit trail is provided in Appendix Sections~\ref{app:dtse_run_manifest}, \ref{app:dtse_scenario_grid}, and~\ref{app:dtse_metric_computation}.

The same principle applies to evidence classes. Mechanism facts are traceable to protocol documentation, and modeled assumptions are traceable to documented parameter ranges and sensitivity handling. Where practitioner context informs a bound, it is marked explicitly as an assumption rather than treated as an empirical estimate. This discipline reduces ambiguity in interpretation and allows later readers to reproduce, audit, or extend the DTSE experiment set without relying on undocumented judgment.

\section{Conclusion}
\label{sec:conclusion}

\subsection{Answer to the Research Question}
This thesis answers RQ1 (Section~\ref{subsec:research_question}) in comparative rather than predictive terms. Within DTSE, DePIN tokenomic mechanisms behave differently under physical and economic stress because their reward rules, sink pathways, and participation conditions create different transmission paths across utilization, provider economics, participation, and incentive--usage alignment.

DTSE therefore shifts the question from ``is a mechanism stable'' to ``which subsystem breaks first under a given channel, and how quickly does that deviation propagate.'' Within the experiment set, the answer appears in first-moving metric families, detection timing, and operational failure-mode signatures (Table~\ref{tab:dtse_cross_scenario_synthesis}; Table~\ref{tab:failure_modes}). Robustness is assessed by where sensitivity concentrates under matched stress inputs---especially in utilization-linked fields, provider-economics and participation fields, and incentive--usage alignment proxies---rather than by point forecasts or prevalence claims about live networks.

For the capped-supply ONO anchor, the results show the main trade-off clearly. User-facing price insulation through Data Credits and fixed-inventory reward logic can reduce one class of volatility transmission, but long-run robustness still depends on whether usage-linked sinks and provider economics remain aligned as stress propagates through the system. More broadly, the thesis shows that neither BME-oriented nor capped-supply designs escape the physical constraints of DePIN. What changes is where the stress appears first and how it accumulates over time.

\subsection{Theoretical and Methodological Contributions}
The thesis contributes a coherent line from stress concepts and mechanism vocabulary to an evaluator that makes assumptions explicit and comparisons auditable. The main contributions are:
\begin{itemize}
    \item \textbf{Methodological framing:} DTSE specifies a context-of-use for simulation-based DePIN stress testing with explicit non-goals and interpretation boundaries.
    \item \textbf{Comparative stress contract:} the thesis defines standardized exogenous stress channels and a fixed metric suite with operational thresholds for cross-profile comparison.
    \item \textbf{Diagnostic vocabulary:} Chapter~8 formalizes failure modes as metric-pattern signatures rather than protocol-level judgments.
    \item \textbf{Interpretive governance layer:} the discussion introduces neutral intervention archetypes while preserving separation between mechanism facts, modeled assumptions, and DTSE outputs.
\end{itemize}

\subsection{Limitations}
The conclusions are bounded by DTSE abstractions and by the scenario and parameter ranges used. Demand regimes are exogenous stylizations, the price-process signal is reduced-form, and provider decision rules simplify participation behavior. Governance interventions, adversarial dynamics, and verification integrity are not evaluated unless they are explicitly modeled and evidenced. DTSE outputs therefore support comparative statements about sensitivity and failure-mode signatures within the DTSE contract (Section~\ref{sec:methodology}); they do not support claims of real-world prevalence, causal certainty, or universal mechanism superiority without additional validation \cite{Morris2019, Braakman2022}.

There is also an evidence-boundary limit. Context sources can describe ecosystem structure and timing during the thesis window, but they are not treated as mechanism-fact proof. Mechanism conclusions remain grounded in protocol or governance documentation and in DTSE artifact outputs. Contextual reporting is used to frame the decision environment and comparator scope, not to replace primary evidence.

\subsection{Future Research Agenda}
\label{subsec:future_research_agenda}
The next extensions follow directly from the evaluator framing and the stated limits:
\begin{enumerate}
    \item Expand mechanism coverage and rule fidelity for additional documented designs while preserving the same stress-channel and metric contracts for comparability.
    \item Strengthen calibration and validation pathways where data is available (e.g., benchmarking selected DTSE outputs against historical stress episodes) without converting DTSE into a forecasting model \cite{McCulloch2022, Collins2024}.
    \item Represent additional decision layers currently out of scope, such as governance-response latency or verification failures, provided evidence boundaries remain explicit and auditable.
    \item Formalize intervention modeling through counterfactual parameter and rule changes evaluated with the same failure-mode vocabulary, preserving separation between modeled intervention effects and claims about real-world feasibility.
\end{enumerate}

\subsection{Closing Statement}
DTSE is presented here as an evaluative instrument for reasoning about DePIN tokenomics under stress when closed-form analysis is insufficient and empirical evidence is incomplete. By making assumptions explicit, reporting outcomes comparatively, and defining failure modes operationally, the framework supports disciplined discussion of robustness without conflating simulation outputs with empirical facts. The thesis therefore contributes both a reproducible structure for stress evaluation and a neutral vocabulary for discussing where DePIN mechanisms fail first, how those failures propagate, and what responses become thinkable once those patterns are visible.
